% Pursuing Stacks (cotes 134-2 to 134-8): archive cote 134-6, site page 122 (PDF page 123), typescript page 440.
% typescript. Draft compared: ps5.tex:7905-7961.
% AI-transcribed by Claude Opus 5.5 (claude-opus-5-5, adaptive thinking, effort high), Message Batches API, batch ps3 (msgbatch_01XdFmuVVBAcgc7EFxdtCqsj), 2026-11-12; not yet checked by eye.
\scanpage{134-6/122}
\mspage{440}

commutative ground ring k,  not only $\underline{Z}$,  with corresponding notion of
ring extension $k\to k'$,  carries over to structures of the type ($\ast$).
Indeed,  for any k-additive category P,  it is easy to define a k'-addi-%
tive category
\[ P\otimes_k k' \quad ,\quad \text{for given homomorphism } k\to k' \ , \]
for instance as the solution of the obvious 2-universal problem \xout{of} \inserted{corresponding to}
mapping P k-additively into k'-additive categories, or more evidently
by taking the same objects as for P,  but with
\[ \mathrm{Hom}_{P'}(a,b) = \mathrm{Hom}_{P}(a,b)\otimes_k k' \quad . \]
Thus any ``coalgebra stucture in (Cat)'' ($\ast$) over the ground ring k,  \xout{give}
gives rise to similar structure over ground ring k'.

Of course,  among the relevant axioms for the diagonal functor ($\ast$),
is the existence of unit objects in $P^{\&}$ and $Q^{\&}$ ,  which may be viewed
equally as \inserted{k-additive} functors (defined up to unique isomorphism)
\[ P^{o}\to Ab_k \quad , \quad P\to Ab_k \qquad , \]
playing the role I would think of ``augmentation'' and ``coaugmentation''
in the more familiar set-up of \inserted{ordinary} coalgebras.  Denoting these objects by
$k_P$  and $k_Q$  respectively (in analogy to the constant presheaves $k_A$  and
$k_B$ on A and B) ,  \xout{Ext} $\mathrm{RHom}(k_P\ ,\ -\ )$  now allows expression of
cohomology or cointegration,  and  $-\overset{L}{\ast}_k(k_Q)$  allows expression of
homology or integration (for complexes in $P^{\&}_{\hand{\mathcal{M}}}$ say).  ``Constant coeffi-%
cients''  on P i.e. in $P^{\&}$ may now be defined,  as objects in $P^{\&}$ of the
type
\[ U\otimes_k k_P \qquad , \]
where U is in $Ab_k$  i.e. is any k-module,  and hence we get homology and
cohomology invariants with coefficients in any such U (or complexes of
such),  and surely too cup and cap products... Also,  quasi-isomorphisms
of structures ($\ast$)  (with units)  can now be defined in an evident way,
hence a derived category which merits to be understood,  when $k=\underline{Z}$ , in
terms of the homotopy category (Hot).  I would'nt expect of course
that for any small category A,  the abelianization Add(A) together with
it's diagonal map allows to recover the fundamental group of A (A 0-con-%
nected say),  still less to recover the homotopy type, unless A is
1-connected.  As was the case visibly for De Rham complexes,  if we
hope to recover general homotopy types (not only 1-connected ones), we
should \xout{look at a} work with slightly more sophisticated structures still,

\margin{122}
